\section{Výzkum}

Tato kapitola se zabývá výzkumem problematiky týkající se samobalančních robotů.\\
Co to vůbec ti samobalanční roboti jsou?\\
Jedná se o typ robota fungující jako otočené kyvadlo. Zjednodušeně se skládá z osy rotace, která se nachází pod jeho těžištěm a 


\subsection{Otočené pendulum}

Otočené pendulum je nestabilní systém, který potřebuje aktivní řízení.
Představte si jej jako misku, kterou otočíte vzhůru nohama a položte na ni kuličku.
V jednom jediném bodě je kulička stabilní, ale když do ní jakkoliv šťouchnete, spadne.
To samé plati i pro otočené kyvadlo, kterým je také tento robot.
\newline

Jak tento problém vyřešit? Co by se toereticky stalo, kdybychom začali touto miskou pohybovat, dokud bychom kuličku opět nedostali do stabilního bodu? O to stejné se pokouší také stabilizace robota. Podle dat naměřených akcelerometrem se tměří náklon, který se poté pohybem robota snaží snížit na nulu (nebo stabilní bod)
\newline

OBRÁZEK ENERGETICKÝCH HLADIN OTOČENÉHO PENDULA

\subsection{metody řízení}
Jaké tedy existují metody měření? Nejčastěji se používá kombinace Arduina, nebo podobného mikrokontroleru s akcelerometrem \ac{MPU6050}. Tato kombinace je jednoduchá pro realizaci díky široké dostupnosti knihoven pro kombinaci Arduino, \ac{MPU6050}. Není však ideální kvůli nízkému výkonu Arduina. Frekvence čipu arduina je $16MHz$. Pro srovnání, frekvence moderního procesoru v počítači je řádové v $GHz$.
\newline

Další rozšířenou alternativou je použití mikrokontroleru \ac{ESP32}. Tento mikrokontroler se dodává v několika variantách. Která z nich je však nejvhodnější?

\subsection{výběr "balení" \ac{ESP32}}
%Výběr formátu
\ac{ESP32} se dá pořídit ve čtyřech verzích. První z nich je \ac{soc}. Je nejmenší, ale nejsložitější pro práci, jelokiž se jedná pouze o mikroprocesor. U některých verzí neobshuje dokonce ani \ac{flash} paměť. Bostrádá také napěťovou regulaci, ochranu vstupů či dokonce anténu, která je pro bezdrátovou komunikaci velmi důležitá.
\newline
% \includegraphics[options]{name}

Přijatelnější verzí je \ac{ESP32} modul. Tento typ balení již obsahuje základní ochranu napájecích vstupů, oscilátor a některé verze dokonce anténu, nebo alespoň konektor pro anténu. Rozměry těchto modulů jsou většinou (19-40)x20mm.
Tato verze obsahuje také paměť \ac{flash} či dokonce paměť \ac{psram}
OBRÁZEK ESP MODULU
\newline

Předposledním způsobem balení je tzv \ac{devboard}. Ta obsahuje modul \ac{ESP32} s dalšími obvody, mezi kterými můžeme najít blink-diodu, USB převodník, či napěťový regulátor na 3V3.
%výběr čipu
TABULKA POROVNÁNÍ ESP32 ČIPŮ



%%% Zdroje pod čarou 
% espressif
% random nerd tutorials
% 